% Part: counterfactuals
% Chapter: introduction
% Section: strict-conditional

\documentclass[../../../include/open-logic-section]{subfiles}

\begin{document}

\olfileid{cnt}{int}{str}

\olsection{严格条件句}

Lewis（刘易斯）引入了\emph{严格条件句}~$\strictif$，并论证它而非实质条件句才对应于蕴涵。
在真势模态逻辑中，$!A \strictif !B$ 可以定义为 $\Box(!A \lif !B)$。 
因此，一个严格条件句（在一个世界上）为真，当且仅当对应的实质条件句是必然的。

严格条件句在应对实质条件句的悖论方面表现如何？ 
前件为假的严格条件句和后件为真的严格条件句，都可能为真，也可能为假。
此外，$(!A \strictif !B) \lor (!B \strictif !A)$ 不是有效的。
严格条件句 $!A \strictif !B$ 也不等价于 $\lnot !A \lor !B$，因此它不是真值函项的。

我们有：
\begin{align}
  !A \strictif !B & \Entails \lnot !A \lor !B
  \text{ 但：}\\
  \lnot !A \lor !B & \Entails/ !A \strictif !B\\
  !B & \Entails/ !A \strictif !B\\
  \lnot !A & \Entails/ !A \strictif !B\\
  \lnot(!A \lif !B) & \Entails/ !A \land \lnot !B
  \text{ 但：}\\
  !A \land \lnot !B & \Entails \lnot(!A \strictif !B)
  \intertext{然而，严格条件句仍然支持肯定前件式：}
  !A, !A \strictif !B & \Entails !B
  \intertext{严格条件句具有聚合性：}
  !A \strictif !B,  !A \strictif !C & \Entails !A \strictif (!B \land !C)
  \intertext{前件强化在严格条件句中成立：}
  !A \strictif !B & \Entails (!A \land !C) \strictif !B
  \intertext{严格条件句也是传递的：}
  !A \strictif !B, !B \strictif !C & \Entails !A \strictif !C
  \intertext{最后，严格条件句等价于它的逆否命题：}
  !A \strictif !B & \Entails \lnot !B \strictif \lnot !A\\
  \lnot !B \strictif \lnot !A & \Entails !A \strictif !B
\end{align}

\begin{prob}
  针对以下对严格条件句不成立的后承关系，给出 \Log{S5} 反例：
  \begin{enumerate}
  \item $\lnot p \Entails/ \Box(p \lif q)$
  \item $q \Entails/ \Box(p \lif q)$
  \item $\lnot\Box(p \lif q) \Entails/ p \land \lnot q$
  \item $\Entails/ \Box(p \lif q) \lor \Box(q \lif p)$
  \end{enumerate}
\end{prob}

\begin{prob}
  通过给出以下公式的 \Log{S5} 证明，说明有效的后承关系对严格条件句成立：
  \begin{enumerate}
  \item  $\Box(!A \lif !B) \Entails \lnot !A \lor !B$
  \item  $!A \land \lnot !B \Entails \lnot\Box(!A \lif !B)$
  \item  $!A, \Box(!A \lif !B) \Entails !B$
  \item  $\Box(!A \lif !B), \Box(!A \lif !C) \Entails \Box(!A \lif (!B
    \land !C))$
  \item  $\Box(!A \lif !B) \Entails \Box((!A \land !C) \lif !B)$
  \item  $\Box(!A \lif !B), \Box(!B \lif !C) \Entails \Box(!A
    \lif !C)$
  \item  $\Box(!A \lif !B) \Entails \Box(\lnot !B \lif \lnot !A)$
  \item  $\Box(\lnot !B \lif \lnot !A) \Entails \Box(!A \lif !B)$
  \end{enumerate}
  \end{prob}

然而，严格条件句仍有其自身的“悖论”。
正如前件为假或后件为真的实质条件句为真，前件\emph{必然}为假或后件必然为真的严格条件句也为真。
此外，任何真的严格条件句都必然为真，任何假的严格条件句都必然为假。
换言之，我们有
\begin{align}
  \Box \lnot !A & \Entails !A \strictif !B\\
  \Box !B & \Entails !A \strictif !B\\
  !A \strictif !B& \Entails \Box(!A \strictif !B)\\
  \lnot(!A \strictif !B) & \Entails \Box\lnot(!A \strictif !B)
\end{align}
如果你将 $\strictif$ 看作“蕴涵”，这些并不是问题。
毕竟，逻辑后承关系是数学事实，因此不可能是偶然的。
但如果你想将 $\strictif$ 用作刻画“如果……那么……”的逻辑联结词，它们确实会引发问题，尤其是最后两条。
因为毫无疑问，存在着偶然为真或偶然为假的“如果……那么……”陈述——事实上，它们通常既非必然也非不可能。

\begin{prob}
  给出以下公式在 \Log{S5} 中的证明：
  \begin{enumerate}
    \item  $\Box \lnot !A \Entails !A \strictif !B$
    \item $!A \strictif !B \Entails \Box(!A \strictif !B)$
    \item $\lnot(!A \strictif !B)  \Entails \Box\lnot(!A \strictif !B)$
  \end{enumerate}
  请利用 $\strictif$ 的定义。
\end{prob}

\end{document}